% ============================================================
%  Section 6 — Conclusion  (placeholder)
% ============================================================

\section{Conclusion}
\label{sec:conclusion}

\noindent\textit{[To be completed. Planned structure:]}

\noindent\textit{\textbf{Summary paragraph.} Restate the main finding: the price-age profile of
top French wines is non-monotone; the trough zone (32--40 years post-vintage) is a structural
feature of demand composition, not noise. Burgundy Grand Cru avoids the trough because its
entry price screens out consumption buyers.}

\noindent\textit{\textbf{Contribution paragraph.} We provide (i) direct evidence on how buyer
heterogeneity shapes asset prices at a key life-cycle transition; (ii) a clean identification
strategy exploiting wine-specific maturity dates; (iii) a tractable structural model that
rationalises the cross-sectional patterns and yields testable comparative statics.}

\noindent\textit{\textbf{Broader implications paragraph.} Results speak to the general class
of storable durable goods with maturity (whisky, aged spirits, vintage cars, art with condition
degradation). The mechanism — exit of consumption buyers at a quality threshold — should appear
in any market with two-type demand and finite consumption value.}

\noindent\textit{\textbf{Limitations and future work.} Endogenous supply (seller timing),
dynamic buyer strategies, and the role of financial intermediaries (wine funds, storage
platforms) are left to future work.}
